\documentclass[11pt]{article}
\usepackage{booktabs,xcolor,colortbl,caption,array,graphicx,amsmath}
\usepackage[landscape,margin=1.2cm]{geometry}
\begin{document}\small\setcounter{table}{6}
\begin{table}[!h]
\centering
\caption{\label{tab:tab:summ_scen}NFA et dark matter sous deux scénarios de taux, moy. 2013--2022 (\% du PIB)}
\centering
\fontsize{10}{12}\selectfont
\begin{tabular}[t]{llrrrrrrr}
\toprule
\multicolumn{2}{c}{\textbf{ }} & \multicolumn{4}{c}{\textbf{Mesures NFA}} & \multicolumn{3}{c}{\textbf{Dark matter}} \\
\cmidrule(l{3pt}r{3pt}){3-6} \cmidrule(l{3pt}r{3pt}){7-9}
Entité & Groupe & NIIP & H\&S (5\%) & NFA$_A$ & NFA$_B$ & DM$_A$ & DM$_B$ & B$-$A\\
\midrule
China & emerging & 14.4 & -9.5 & NaN & NaN & NaN & NaN & NaN\\
European Union & advanced & -8.5 & 0.5 & -8.8 & -8.8 & -0.3 & -0.3 & 0.0\\
Japan & advanced & 63.0 & 81.6 & 73.5 & 73.5 & 10.5 & 10.5 & 0.0\\
United States & privilege & -52.7 & 19.7 & -8.8 & -26.4 & 44.0 & 26.4 & -17.6\\
\bottomrule
\multicolumn{9}{l}{\rule{0pt}{1em}\textit{\textit{Notes:} }}\\
\multicolumn{9}{l}{\rule{0pt}{1em}Scén. A: taux G\&R (2006) universels. Scén. B: Gourinchas, Rey \& Govillot (2017). Privilège (USA, CHE, GBR): r\textsubscript{FDI}=9\%, r\textsubscript{dette}=2\%. Émergents (CHN): r\textsubscript{FDI}=6\%, r\textsubscript{dette}=5\%. Un r plus \'{e}lev\'{e} implique un stock capitalis\'{e} plus faible pour le m\^{e}me flux de revenu.}\\
\end{tabular}
\end{table}
\end{document}
